

\clearpage
\section{Presupuesto y financiamiento}
En este apartado se deben detallar los gastos que implican cada una de las acciones o actividades del proyecto, así como los materiales que serán empleados. Se recomienda usar una tabla de resumen de costos y algunas tablas de detalle. Los rubros del presupuesto podrán ser: a. Visitas de campo, b. Materiales, c. Insumos, d. Pago de servicios, e. Capacitación, f. Difusión de resultados.
\subsection{Talento humano}
En este apartado se deben detallar los gastos que implican cada uno de los miembros del proyecto principalmente el estudiante y el director de trabajo de titulación. Se debe calcular e costo por hora que se va a utilizar y multiplicarlo por el total de horas Tabla \begin{table}[H]
\centering
\caption{Presupuesto Talento Humano}
\label{tab:presupuesto-talento}
\begin{tabular}{|l|c|c|c|c|}
\hline
\rowcolor[HTML]{B6DAF8}
\textbf{Descripción} & \textbf{Cantidad} & \textbf{Unidad} & \textbf{Valor Unitario} & \textbf{Valor total} \\ \hline
Estudiante           & 360               & hora            & \$10                    & \$3.600,00            \\ \hline
Director de TIC      & 16                & hora            & \$10.48                 & \$167,68              \\ \hline
\multicolumn{4}{|r|}{\textbf{Total}}                         & \textbf{\$3.767,68}     \\ \hline
\end{tabular}
\end{table}


\subsection{Servicios}
En este apartado se deben detallar los gastos que implican cada uno de los miembros del proyecto principalmente el estudiante y el director de trabajo de titulación. Se debe calcular e costo por hora que se va a utilizar y multiplicarlo por el total de horas Tabla \ref{tab:presupuesto-servicios}.
\begin{table}[H]
\centering
\caption{Presupuesto Servicios}
\label{tab:presupuesto-servicios}
\begin{tabular}{|l|c|c|c|c|}
\hline
\rowcolor[HTML]{B6DAF8}
\textbf{Descripción} & \textbf{Cantidad} & \textbf{Unidad} & \textbf{Valor Unitario} & \textbf{Valor total} \\ \hline
Internet             & 4                & Mensualidades    & \$22                     & \$88,00               \\ \hline
Firma electrónica    & 1                & Licencia         & \$20                     & \$20,00               \\ \hline
\multicolumn{4}{|r|}{\textbf{Total}}                          & \textbf{\$108,00}        \\ \hline
\end{tabular}
\end{table}

\subsection{Hardware}
En este apartado se deben detallar los gastos que implican cada uno de los miembros del proyecto principalmente el estudiante y el director de trabajo de titulación. Se debe calcular e costo por hora que se va a utilizar y multiplicarlo por el total de horas Tabla \ref{tab:presupuesto-hardware}.
\begin{table}[H]
\centering
\caption{Presupuesto Hardware}
\label{tab:presupuesto-hardware}
\begin{tabular}{|l|c|l|c|c|}
\hline
\rowcolor[HTML]{B6DAF8}
\textbf{Descripción} & \textbf{Cantidad} & \textbf{Unidad} & \textbf{Valor Unitario} & \textbf{Valor total} \\ \hline
Laptop               & 1                & Computadora     & \$22                    & \$88,00              \\ \hline
Jetson Nano          & 1                & Placa           & \$20                    & \$20,00              \\ \hline
\multicolumn{4}{|r|}{\textbf{Total}}                         & \textbf{\$108,00}       \\ \hline
\end{tabular}
\end{table}

\subsection{Software}
En este apartado se deben detallar los gastos que implican cada uno de los miembros del proyecto principalmente el estudiante y el director de trabajo de titulación. Se debe calcular e costo por hora que se va a utilizar y multiplicarlo por el total de horas Tabla \ref{tab:presupuesto-software}.

\begin{table}[H]
\centering
\caption{Presupuesto Software}
\label{tab:presupuesto-software}
\begin{tabular}{|l|c|l|c|c|}
\hline
\rowcolor[HTML]{B6DAF8}
\textbf{Descripción} & \textbf{Cantidad} & \textbf{Unidad} & \textbf{Valor Unitario} & \textbf{Valor total} \\ \hline
Google Colab         & 4                & Licencias       & \$0                     & \$00,00              \\ \hline
SPSS                 & 4                & Licencias       & \$20                    & \$80,00              \\ \hline
\multicolumn{4}{|r|}{\textbf{Total}}                         & \textbf{\$80,00}        \\ \hline
\end{tabular}
\end{table}

\subsection{Resumen del presupuesto}
En este apartado se deben detallar los gastos que implican cada uno de los miembros del proyecto principalmente el estudiante y el director de trabajo de titulación. Se debe calcular e costo por hora que se va a utilizar y multiplicarlo por el total de horas Tabla \ref{tab:resumen-presupuesto}

\begin{table}[H]
\centering
\caption{Resumen del presupuesto}
\label{tab:resumen-presupuesto}
\begin{tabular}{|l|c|}
\hline
\rowcolor[HTML]{B6DAF8}
\textbf{Descripción}     & \textbf{Valor}    \\ \hline
Talento Humano           & \$3.767,68        \\ \hline
Servicios                & \$108,00          \\ \hline
Hardware                 & \$108,00          \\ \hline
Software                 & \$80,00           \\ \hline
\textbf{Subtotal}        & \textbf{\$4.063,68} \\ \hline
Imprevisto 10\%          & 406,63            \\ \hline
\rowcolor[HTML]{5BC0DE}
\textbf{Total}           & \textbf{\$4.470,31} \\ \hline
\end{tabular}
\end{table}
